\documentclass{exam-zh}
\usepackage{edu-practice}

\examsetup{
  question/show-answer = false,
  fillin/show-answer = false,
  solution/show-solution = hide,
}

% ---------- 教师版解答样式 ----------
\newtcolorbox{solutionblock}{
  enhanced, breakable,
  colback=edu-blue-bg,
  colframe=edu-blue!25,
  boxrule=0pt,
  borderline west={2pt}{0pt}{edu-blue!50},
  arc=0pt,
  left=3mm, right=3mm, top=2mm, bottom=2mm,
  before skip=2mm,
  after skip=3mm,
  fontupper=\small,
}

% 解答步骤编号
\newcounter{solstep}
\newcommand{\solstep}[2]{%
  \stepcounter{solstep}%
  \par\medskip
  \noindent\hangindent=6mm
  {\color{edu-blue}\bfseries\textcircled{\scriptsize\thesolstep}\enspace #1}\enspace #2\par
}

\begin{document}

\section*{立体几何距离与角 · 练习}

\section*{一、选择题（每题 8 分，共 16 分）}


\begin{question}[points=8]
  直线 $l$ 与平面 $\alpha$ 斜交，$l$ 上一点 $P$（斜足外）在 $\alpha$ 内的射影为 $P'$，斜足为 $A$。下列说法正确的是（　　）。
  \begin{hintbox}
\textbf{提示：}线面角 = 斜线与其在面内射影的夹角。先有垂线才有射影。\end{hintbox}

  \begin{choices}
    \item 直线 $l$ 与平面 $\alpha$ 所成的角是 $\angle PAP'$
    \item 直线 $l$ 与平面 $\alpha$ 所成的角是 $\angle lAA'$，其中 $AA'\perp\alpha$
    \item 直线 $l$ 与平面 $\alpha$ 所成的角是 $l$ 与 $\alpha$ 内任意一条直线的夹角
    \item 线面角可以大于 $90^\circ$
  \end{choices}
\end{question}



\begin{question}[points=8]
  用余弦定理算得两条异面直线平移后的夹角 $\theta$ 满足 $\cos\theta=-\dfrac{\sqrt2}{2}$。则这两条异面直线所成角为（　　）。
  \begin{hintbox}
\textbf{提示：}异面直线所成角范围是 $(0^\circ,90^\circ]$，算出钝角要取补角。\end{hintbox}

  \begin{choices}
    \item $45^\circ$
    \item $135^\circ$
    \item $60^\circ$
    \item $120^\circ$
  \end{choices}
\end{question}


\section*{二、填空题（每题 10 分，共 20 分）}


\begin{question}[points=10]
  正方体 $ABCD\text{-}A_1B_1C_1D_1$ 的棱长为 $2$。则顶点 $A_1$ 到对角面 $BDD_1B_1$ 的距离为 $\underline{\hspace{2em}}$。
  \begin{hintbox}
\textbf{提示：}点面距、垂足位置不显然时，用等体积法：选一个三棱锥列 $V$ 相等。可试 $V_{A_1\text{-}BDD_1}=V_{D\text{-}A_1BD_1}$。\end{hintbox}

\end{question}



\begin{question}[points=10]
  正四棱锥 $P\text{-}ABCD$ 的底面边长为 $2$，侧棱 $PA=3$，$O$ 为底面中心。则侧棱 $PA$ 与底面 $ABCD$ 所成角的正切值为 $\underline{\hspace{2em}}$。
  \begin{hintbox}
\textbf{提示：}侧棱与底面的线面角 = 斜线 $PA$ 与射影 $AO$ 的夹角。先求高 $PO$ 和 $AO$。\end{hintbox}

\end{question}


\section*{三、解答题（共 60 分）}

\needspace{8\baselineskip}

\begin{problem}[points=12]
    正方体 $ABCD\text{-}A_1B_1C_1D_1$ 的棱长为 $2$，$O$ 为 $AC$ 与 $BD$ 的交点。
\begin{enumerate}[label=(\arabic*)]
  \item 证明 $BD_1\perp AC$；
  \item 求截面 $AB_1C$ 与底面 $ABCD$ 所成二面角 $A\text{-}AC\text{-}B_1$ 的余弦值。
\end{enumerate}

    \begin{hintbox}
\textbf{提示：}（1）用三垂线定理：$BD\perp AC$（正方形对角线互相垂直），$D_1D\perp$ 底面，得 $BD_1\perp AC$。（2）截面 $AB_1C$ 与底面的棱为 $AC$，取 $AC$ 中点 $O$，$OB\perp AC$，由三垂线 $OB_1\perp AC$，$\angle B_1OB$ 是平面角。\end{hintbox}

  \answerarea[62mm]
\end{problem}


\needspace{8\baselineskip}

\begin{problem}[points=14]
    在三棱锥 $P\text{-}ABC$ 中，$PA\perp$ 平面 $ABC$，$PA=2$，$AB=AC=2$，$\angle BAC=120^{\circ}$。求点 $A$ 到平面 $PBC$ 的距离。

    \begin{hintbox}
\textbf{提示：}点 $A$ 到面 $PBC$ 的距离用等体积法：$V_{P\text{-}ABC}=V_{A\text{-}PBC}$。先算 $S_{\triangle ABC}$ 和 $S_{\triangle PBC}$。注意 $PA\perp$ 底面，可先算 $PB$、$PC$。\end{hintbox}

  \answerarea[70mm]
\end{problem}


\needspace{8\baselineskip}

\begin{problem}[points=12]
    在正三棱柱 $ABC\text{-}A_1B_1C_1$ 中，底面边长为 $2$，侧棱 $AA_1=2$，$D$ 为 $BC$ 的中点。
\begin{enumerate}[label=(\arabic*)]
  \item 证明 $A_1D\perp BC$；
  \item 求直线 $A_1B$ 与底面 $ABC$ 所成角的正切值。
\end{enumerate}

    \begin{hintbox}
\textbf{提示：}（1）$AD\perp BC$（正三角形中线即高），$A_1A\perp$ 底面，由三垂线定理得 $A_1D\perp BC$。（2）$A_1B$ 与底面所成角 = $A_1B$ 与射影 $AB$ 的夹角，在 $\text{Rt}\triangle A_1AB$ 中算。\end{hintbox}

  \answerarea[62mm]
\end{problem}


\needspace{8\baselineskip}

\begin{problem}[points=10]
    正方体 $ABCD\text{-}A_1B_1C_1D_1$ 的棱长为 $2$。求点 $D_1$ 到直线 $AC$ 的距离。

    \begin{hintbox}
\textbf{提示：}连 $BD$ 交 $AC$ 于 $O$。$D_1D\perp$ 底面，$BD\perp AC$，由三垂线定理 $D_1O\perp AC$，$D_1O$ 即为所求。\end{hintbox}

  \answerarea[55mm]
\end{problem}


\needspace{8\baselineskip}

\begin{problem}[points=12]
    在棱长为 $2$ 的正四面体 $P\text{-}ABC$ 中，$O$ 为底面 $\triangle ABC$ 的中心。
\begin{enumerate}[label=(\arabic*)]
  \item 求点 $P$ 到平面 $ABC$ 的距离；
  \item 设 $E$、$F$ 分别为 $PA$、$BC$ 的中点，求异面直线 $PA$ 与 $BC$ 的距离。
\end{enumerate}

    \begin{hintbox}
\textbf{提示：}（1）$PO\perp$ 底面，在 $\text{Rt}\triangle PAM$（$M$ 为 $BC$ 中点）里算 $PO$。（2）$E$、$F$ 为 $PA$、$BC$ 中点，$EF$ 为异面直线 $PA$ 与 $BC$ 的公垂线段，在 $\triangle PAM$ 里求 $EF$。\end{hintbox}

  \answerarea[70mm]
\end{problem}



\end{document}
